\documentclass[aspectratio=169,10pt]{beamer}

\usetheme{Madrid}
\usepackage[T1]{fontenc}
\usepackage[utf8]{inputenc}
\usepackage{lmodern}
\usepackage{amsmath,amssymb}
\usepackage{booktabs}
\usepackage{listings}
\usepackage{xcolor}

\definecolor{WorkshopRed}{RGB}{166,20,45}
\definecolor{CodeBackground}{RGB}{246,246,246}
\definecolor{CodeComment}{RGB}{90,110,90}

\setbeamercolor{structure}{fg=WorkshopRed}
\setbeamercolor{alerted text}{fg=WorkshopRed}
\setbeamercolor{block title}{fg=white,bg=WorkshopRed}
\setbeamercolor{block body}{bg=WorkshopRed!5}
\setbeamertemplate{navigation symbols}{}
\setbeamertemplate{footline}{}

\lstset{
  basicstyle=\ttfamily\scriptsize,
  backgroundcolor=\color{CodeBackground},
  commentstyle=\color{CodeComment},
  keywordstyle=\color{WorkshopRed}\bfseries,
  frame=single,
  framerule=0.3pt,
  columns=fullflexible,
  keepspaces=true,
  showstringspaces=false,
  breaklines=true,
  morekeywords={def,theorem,structure,inductive,where,by,match,with,if,then,else,forall,True,False}
}

\title{From Hallucination to Proof}
\subtitle{Lean 4 and software verification for AI-assisted robotics}
\author{Łukasz Janiec}
\institute{Wrocław University of Science and Technology}
\date{}

\begin{document}

\begin{frame}
  \titlepage
  \begin{center}
    \alert{AI can propose. The verifier decides what is accepted.}
  \end{center}
\end{frame}

\begin{frame}{Why verification matters now}
  \begin{itemize}
    \item AI agents can generate, refactor, and explain software very quickly.
    \item Robotics software combines discrete state, asynchronous events,
          sensors, resources, communication, and physical effects.
    \item Plausible code is useful, but plausibility is not a correctness
          certificate.
    \item We need a boundary between \emph{producing a candidate} and
          \emph{accepting it as correct}.
  \end{itemize}

  \begin{block}{Workshop principle}
    The model may be wrong. The checker must remain independent of the model.
  \end{block}
\end{frame}

\begin{frame}{Testing and proving answer different questions}
  \begin{columns}[T,onlytextwidth]
    \column{0.47\textwidth}
      \begin{block}{Testing}
        \begin{itemize}
          \item Executes selected examples.
          \item Finds concrete failures.
          \item Scales well to real implementations.
          \item Cannot normally cover every state and trace.
        \end{itemize}
      \end{block}

    \column{0.47\textwidth}
      \begin{block}{Formal proof}
        \begin{itemize}
          \item States an explicit mathematical claim.
          \item Covers every case represented by the model.
          \item Makes assumptions visible.
          \item Is only as relevant as the model and specification.
        \end{itemize}
      \end{block}
  \end{columns}

  \vspace{0.5em}
  \centering
  \alert{Python shows what happens. Lean checks what must always happen.}
\end{frame}

\begin{frame}[fragile]{What Lean 4 checks}
  \begin{lstlisting}
structure State where
  active : Bool
  sensorOpen : Bool

def Safe (s : State) : Prop :=
  s.sensorOpen = true -> s.active = true

theorem step_preserves_safe
    (s : State) (e : Event)
    (h : Safe s) :
    Safe (step s e) := by
  -- proof checked by Lean
  \end{lstlisting}

  \begin{itemize}
    \item A proposition is represented as a type.
    \item A proof is a term inhabiting that type.
    \item Tactics construct proof terms; they are not the final authority.
    \item Lean's kernel checks the resulting term against the theorem statement.
  \end{itemize}
\end{frame}

\begin{frame}{Why the kernel matters}
  \begin{center}
    \Large
    Human or AI proposal
    $\longrightarrow$
    proof term
    $\longrightarrow$
    \alert{Lean kernel}
    $\longrightarrow$
    accepted or rejected
  \end{center}

  \vspace{1em}
  \begin{itemize}
    \item Gemini, Codex, Claude, Lean Copilot, or a local model may suggest a proof.
    \item The suggestion may contain a syntax error, a nonexistent lemma, or a
          logically invalid step.
    \item The trusted result is the kernel-checked theorem, not the confidence
          or reputation of the generator.
  \end{itemize}

  \begin{block}{Useful architecture}
    Strong, fallible generator outside a small, strict verification boundary.
  \end{block}
\end{frame}

\begin{frame}{What have we actually verified?}
  \begin{center}
    \begin{tabular}{c}
      Real robot and environment \\
      $\downarrow$ modelling assumptions \\
      Formal state and transition model \\
      $\downarrow$ specification \\
      Theorem: invariant, refinement, or contract \\
      $\downarrow$ proof \\
      \alert{Kernel-checked result}
    \end{tabular}
  \end{center}

  \vspace{0.7em}
  \begin{enumerate}
    \item Is the theorem true of the formal model?
    \item Does the model capture the relevant part of the real system?
    \item Did we formalize the right requirement?
  \end{enumerate}

  Lean answers the first question extremely well. The other two remain
  engineering and specification work.
\end{frame}

\begin{frame}{Example 1: lifecycle state is not the whole state}
  \begin{columns}[T,onlytextwidth]
    \column{0.50\textwidth}
      \textbf{Modeled application state}
      \begin{itemize}
        \item lifecycle mode;
        \item sensor handle;
        \item log resource;
        \item actuator authority.
      \end{itemize}

      \textbf{Invariant}
      \[
        \mathrm{Active} \Rightarrow
        \mathrm{sensorOpen} \land \mathrm{logOpen}
      \]
      \[
        \mathrm{sensorOpen} \lor \mathrm{logOpen} \lor
        \mathrm{actuatorEnabled}
        \Rightarrow \mathrm{Active}
      \]

    \column{0.46\textwidth}
      \begin{alertblock}{Bug}
        A deactivation changes only the lifecycle enum and leaves the sensor,
        log, or actuator state active.
      \end{alertblock}

      \begin{block}{Formal result}
        Prove once that every modeled transition preserves the invariant, then
        lift that result to all modeled executions.
      \end{block}
  \end{columns}
\end{frame}

\begin{frame}{Example 2: did the refactor preserve behavior?}
  \textbf{Reference Behavior Tree priority}
  \begin{enumerate}
    \item cancel requested $\rightarrow$ stop and cancel;
    \item goal reached $\rightarrow$ report success;
    \item path available $\rightarrow$ drive;
    \item otherwise $\rightarrow$ request a plan.
  \end{enumerate}

  \begin{alertblock}{Regression in the handwritten FSM}
    The rewrite checks \texttt{pathAvailable} before \texttt{cancelRequested}.
    When both are true, the BT stops, while the FSM drives.
  \end{alertblock}

  \begin{block}{Refinement claim}
    For every modeled observation, the corrected FSM returns the same observable
    decision as the reference policy.
  \end{block}
\end{frame}

\begin{frame}{AI-assisted verification loop}
  \begin{enumerate}
    \item The agent reads the goal and proposes a candidate.
    \item The repository runs the exact local checker.
    \item Lean rejects an invalid candidate and returns a concrete diagnostic.
    \item The diagnostic is given back to the agent.
    \item The agent proposes a minimal repair.
    \item The same checker runs again.
  \end{enumerate}

  \vspace{0.5em}
  \begin{center}
    \Large
    proposal $\rightarrow$ rejection $\rightarrow$ feedback
    $\rightarrow$ repair $\rightarrow$ verification
  \end{center}

  \begin{block}{Important distinction}
    A bad proof of a good theorem can be rejected. A good proof of the wrong
    theorem may still compile.
  \end{block}
\end{frame}

\begin{frame}{What formal verification does---and does not---give us}
  \begin{itemize}
    \item It gives strong guarantees for explicit claims under explicit
          definitions and assumptions.
    \item It can expose missing cases, regressions, inconsistent state, and
          false universal claims.
    \item It does not automatically validate sensor models, timing assumptions,
          physical dynamics, or informal requirements.
    \item A formally verified answer to the wrong question is still the wrong
          engineering result.
  \end{itemize}

  \begin{block}{Goal of the workshop}
    Build small, explicit, independently checked contracts at important
    software boundaries---then let AI help without letting it certify itself.
  \end{block}

  \begin{center}
    \alert{Now we will apply this pattern to the workshop tasks.}
  \end{center}
\end{frame}

\end{document}
